% Étude de cas : Sommeil à pattern narcoleptique sur polysomnographie
% Patiente : aeiuno
% Créé : 2026-03-22
% Confidentialité : Anonymisé avec consentement de la patiente
% Statut : Bilan complémentaire proposé — à discuter avec neurologue/somnologue

\documentclass[11pt,a4paper]{article}

\usepackage[utf8]{inputenc}
\usepackage[T1]{fontenc}
\usepackage[french]{babel}
\usepackage{newpxtext}         % Palatino (texte)
\usepackage{newpxmath}         % Palatino (maths)
\usepackage{booktabs}
\usepackage{hyperref}
\usepackage[table]{xcolor}
\usepackage{tcolorbox}
\tcbuselibrary{breakable}
\usepackage{enumitem}
\usepackage{geometry}
\usepackage{float}
\usepackage{tabularx}
\geometry{margin=2.5cm}

% Color definitions
\definecolor{keyfindingbg}{RGB}{212,237,218}
\definecolor{keyfindingborder}{RGB}{40,167,69}
\definecolor{hypothesisbg}{RGB}{217,237,247}
\definecolor{hypothesisborder}{RGB}{23,162,184}
\definecolor{cautionbg}{RGB}{248,215,218}
\definecolor{cautionborder}{RGB}{220,53,69}
\definecolor{mechanismbg}{RGB}{232,232,232}
\definecolor{mechanismborder}{RGB}{108,117,125}
\definecolor{questionbg}{RGB}{255,243,205}
\definecolor{questionborder}{RGB}{255,193,7}

\newtcolorbox{keyfinding}[1][]{
  colback=keyfindingbg, colframe=keyfindingborder,
  adjusted title={\textbf{Constat clé~: #1}}, fonttitle=\bfseries, sharp corners, boxrule=1pt,
  breakable
}

\newtcolorbox{hypothesis}[1][]{
  colback=hypothesisbg, colframe=hypothesisborder,
  adjusted title={\textbf{Hypothèse~: #1}}, fonttitle=\bfseries, sharp corners, boxrule=1pt,
  breakable
}

\newtcolorbox{caution}[1][]{
  colback=cautionbg, colframe=cautionborder,
  adjusted title={\textbf{Attention~: #1}}, fonttitle=\bfseries, sharp corners, boxrule=1pt,
  breakable
}

\newtcolorbox{mechanism}[1][]{
  colback=mechanismbg, colframe=mechanismborder,
  adjusted title={\textbf{#1}}, fonttitle=\bfseries, sharp corners, boxrule=1pt,
  breakable
}

\newtcolorbox{question}[1][]{
  colback=questionbg, colframe=questionborder,
  adjusted title={\textbf{À clarifier~: #1}}, fonttitle=\bfseries, sharp corners, boxrule=1pt,
  breakable
}

\title{Étude de cas~: Sommeil à pattern narcoleptique\\[0.5em]
\large Aeiuno --- Polysomnographie évocatrice de narcolepsie\\[0.3em]
\normalsize Contexte d'EM/SFC --- Bilan complémentaire proposé}
\author{Document de travail --- À discuter avec un neurologue / somnologue}
\date{22 mars 2026}

\begin{document}

\maketitle

\begin{caution}{Document préliminaire}
Cette étude de cas repose sur des données \textbf{limitées}. Les recommandations ci-dessous orientent la démarche diagnostique et thérapeutique et ne doivent pas être interprétées comme des conclusions définitives. \textbf{Toutes les recommandations nécessitent une validation par le médecin traitant.}
\end{caution}

\tableofcontents
\newpage

%=============================================================================
\section{Contexte clinique}
\label{sec:contexte}
%=============================================================================

\begin{table}[H]
\centering
\begin{tabularx}{\textwidth}{l>{\raggedright\arraybackslash}X}
\toprule
\textbf{Paramètre} & \textbf{Valeur} \\
\midrule
Pseudonyme & Aeiuno \\
Sexe & Femme \\
Date de naissance & 8 septembre 1988 \\
Âge & 37 ans \\
Diagnostic principal & Encéphalomyélite myalgique / Syndrome de fatigue chronique (EM/SFC) \\
Examen en question & Polysomnographie (PSG) \\
Résultat & Pattern de sommeil évocateur de narcolepsie \\
\bottomrule
\end{tabularx}
\caption{Résumé clinique}
\end{table}

%=============================================================================
\section{Constat polysomnographique}
\label{sec:constat}
%=============================================================================

\begin{keyfinding}{Pattern narcoleptique à la PSG}
La polysomnographie montre des anomalies compatibles avec un diagnostic de narcolepsie~: latence de sommeil paradoxal (REM) raccourcie et/ou endormissements en sommeil paradoxal (SOREMP). Ces anomalies, chez une patiente atteinte d'EM/SFC, posent la question d'une narcolepsie de type~1 comorbide versus une suppression fonctionnelle du système orexinergique liée à la neuroinflammation de l'EM/SFC.
\end{keyfinding}

%=============================================================================
\section{Diagnostic différentiel~: narcolepsie vraie vs.\ suppression fonctionnelle}
\label{sec:differentiel}
%=============================================================================

Le pattern narcoleptique à la PSG peut refléter deux mécanismes distincts avec des implications thérapeutiques différentes~:

\begin{hypothesis}{Narcolepsie de type~1 comorbide}
Destruction auto-immune des neurones à hypocrétine/orexine dans l'hypothalamus latéral. Shan et al.\ (2026, \textit{Annals of Neurology}) ont confirmé que les lymphocytes T~CD4+ ciblent sélectivement les neurones producteurs d'hypocrétine, avec une densité 11 fois supérieure aux autres types de lymphocytes~T dans la zone atteinte. La narcolepsie de type~1 est fréquemment post-infectieuse (H1N1, streptocoque) et fortement associée au HLA-DQB1*06:02.

\textbf{Arguments pour~:}
\begin{itemize}[nosep]
    \item Pattern PSG spécifique (SOREMPs, latence REM courte)
    \item Contexte post-infectieux compatible (l'EM/SFC et la narcolepsie partagent ce mode de déclenchement)
    \item Cas documenté de comorbidité EM/SFC + POTS + narcolepsie (Liao et al.\ 2021)
\end{itemize}
\end{hypothesis}

\begin{hypothesis}{Suppression fonctionnelle du système orexinergique}
Inhibition des neurones à orexine par les cytokines inflammatoires, sans destruction neuronale. La PGE\textsubscript{2} produite lors de la neuroinflammation supprime l'activité des neurones orexinergiques via les récepteurs EP3/EP4. Des taux réduits d'orexine-A ont été rapportés dans l'EM/SFC (López-Amador 2025, revue systématique de 27 études).

\textbf{Arguments pour~:}
\begin{itemize}[nosep]
    \item Contexte d'EM/SFC avec neuroinflammation documentée
    \item Taux d'orexine-A réduits mais non effondrés dans l'EM/SFC
    \item Mécanisme réversible (contrairement à la destruction neuronale)
\end{itemize}
\end{hypothesis}

%=============================================================================
\section{Bilan complémentaire recommandé}
\label{sec:bilan}
%=============================================================================

\begin{mechanism}{Examens pour trancher entre les deux hypothèses}

\begin{enumerate}
    \item \textbf{Hypocrétine-1 dans le LCR} --- Examen de référence
    \begin{itemize}[nosep]
        \item $< 110$~pg/mL $=$ narcolepsie de type~1 confirmée (critère ICSD-3 suffisant à lui seul)
        \item Normal ou bas-normal $=$ suppression fonctionnelle probable
    \end{itemize}

    \item \textbf{Typage HLA-DQB1*06:02}
    \begin{itemize}[nosep]
        \item Présent chez $\sim$98\% des narcolepsies de type~1
        \item Si négatif~: narcolepsie de type~1 très peu probable
        \item Examen simple (prise de sang), à réaliser en premier
    \end{itemize}

    \item \textbf{Test itératif de latence d'endormissement (TILE / MSLT)} si non réalisé
    \begin{itemize}[nosep]
        \item Critères narcolepsie~: latence moyenne $\leq 8$~min et $\geq 2$~SOREMPs
        \item Un SOREMP dans les 15~min sur la PSG nocturne précédente peut remplacer un SOREMP du TILE
    \end{itemize}
\end{enumerate}

\textbf{Critères ICSD-3 pour la narcolepsie de type~1~:}
\begin{enumerate}[nosep]
    \item Somnolence diurne excessive quotidienne depuis $\geq 3$~mois
    \item Un ou les deux critères suivants~:
    \begin{itemize}[nosep]
        \item Cataplexie + TILE positif (latence $\leq 8$~min, $\geq 2$~SOREMPs)
        \item Hypocrétine-1 dans le LCR $< 110$~pg/mL (ou $< 1/3$ des valeurs normales moyennes)
    \end{itemize}
\end{enumerate}
Le critère LCR est suffisant à lui seul~: la cataplexie et le TILE ne sont pas nécessaires si l'hypocrétine est basse.
\end{mechanism}

%=============================================================================
\section{Stratégie thérapeutique selon le diagnostic}
\label{sec:traitement}
%=============================================================================

\subsection{Si narcolepsie de type~1 confirmée}

\begin{table}[H]
\centering
\begin{tabularx}{\textwidth}{lX}
\toprule
\textbf{Indication} & \textbf{Traitement} \\
\midrule
Somnolence diurne & Modafinil 100--200~mg/j (active les neurones à orexine~; pas de tolérance documentée jusqu'à 40 semaines) \\
Consolidation du sommeil & Oxybate de sodium (Xyrem\textsuperscript{\textregistered}) --- traitement de référence de la narcolepsie \\
Immunothérapie précoce & À discuter~: IgIV, étant donné la confirmation récente du mécanisme auto-immun. Pourrait prévenir la destruction neuronale si prise en charge précoce \\
\midrule
\textbf{Contre-indication} & \textbf{Antagonistes des récepteurs de l'orexine (DORA)} --- suvorexant, lemborexant, daridorexant. Bloquer un système orexinergique déjà détruit est contre-indiqué \\
\bottomrule
\end{tabularx}
\caption{Traitement si narcolepsie de type~1}
\end{table}

\subsection{Si suppression fonctionnelle (EM/SFC)}

\begin{table}[H]
\centering
\begin{tabularx}{\textwidth}{lX}
\toprule
\textbf{Indication} & \textbf{Traitement} \\
\midrule
Somnolence diurne & Modafinil 100--200~mg/j (même justification~: activation des neurones à orexine) \\
Consolidation du sommeil & Daridorexant 25--50~mg --- les DORA sont appropriés ici car les neurones sont supprimés, pas détruits \\
Mécanisme sous-jacent & Cibler la neuroinflammation responsable de la suppression orexinergique \\
\bottomrule
\end{tabularx}
\caption{Traitement si suppression fonctionnelle}
\end{table}

\subsection{Dans les deux cas}

\begin{keyfinding}{Premier pas pragmatique}
Le modafinil est le traitement de première ligne quel que soit le diagnostic final. Commencer à 100~mg, titrer progressivement. \textbf{Surveiller l'apparition d'un malaise post-effort (MPE/PEM)}~: tous les psychostimulants comportent ce risque dans le contexte de l'EM/SFC.
\end{keyfinding}

%=============================================================================
\section{Contexte scientifique récent}
\label{sec:contexte-scientifique}
%=============================================================================

\begin{mechanism}{Narcolepsie confirmée auto-immune (Shan et al.\ 2026)}
Shan et al.\ (\textit{Annals of Neurology}, mars 2026) ont analysé du tissu cérébral post-mortem via la Netherlands Brain Bank et confirmé que la narcolepsie de type~1 résulte d'une réponse auto-immune chronique~:

\begin{itemize}[nosep]
    \item Les lymphocytes T~CD4+ sont 11 fois plus nombreux dans la région des neurones à hypocrétine que les autres types de lymphocytes~T
    \item Ces cellules immunitaires ciblent sélectivement et détruisent les neurones producteurs d'hypocrétine/orexine
    \item Les signatures immunitaires sont détectables des \textbf{décennies} après la destruction initiale
    \item Implication~: une immunothérapie préventive pourrait stopper la perte neuronale si mise en place précocement
\end{itemize}
\end{mechanism}

\begin{mechanism}{Pertinence pour l'EM/SFC}
\begin{itemize}[nosep]
    \item \textbf{Système orexinergique~:} Des taux réduits d'orexine-A sont rapportés dans l'EM/SFC (López-Amador 2025, 27 études)
    \item \textbf{Déclencheur auto-immun post-infectieux~:} Les deux pathologies surviennent après infection via attaque auto-immune sur des cibles du SNC
    \item \textbf{Comorbidité documentée~:} Cas d'EM/SFC + POTS + narcolepsie avec mutations MTHFR (Liao et al.\ 2021)
    \item \textbf{Chevauchement thérapeutique~:} L'immunomodulation (IgIV, rituximab) est étudiée dans les deux conditions
\end{itemize}
\end{mechanism}

%=============================================================================
\section{Références}
\label{sec:references}
%=============================================================================

\begin{enumerate}[nosep]
    \item Shan L et al.\ (2026). Elevated CD4+ T-cells in the hypocretin-producing region confirm autoimmune pathogenesis in type-1 narcolepsy. \textit{Annals of Neurology}.
    \item López-Amador N et al.\ (2025). Orexin System Dysfunction in ME/CFS. Revue systématique de 27 études.
    \item Liao Y et al.\ (2021). Comorbidity of CFS, POTS, and narcolepsy with MTHFR mutation. \textit{Chinese Medical Journal}, 134(12):1495--1497.
    \item American Academy of Sleep Medicine (2014). International Classification of Sleep Disorders, 3rd ed.\ (ICSD-3). Critères diagnostiques de la narcolepsie de type~1.
\end{enumerate}

\end{document}
